\section{Smooth transverse moments and collision-stable dynamics}
\label{sec:smooth-transverse}

The next observable changes the statistical object. For a nonnegative
\(a\in C_c^2(\R)\), define the weighted one-point transverse measure
and its second moment by
\begin{equation}\label{eq:transverse-measure}
 \mu_{\lambda,a}=\sum_j a(x_j(\lambda))\delta_{y_j(\lambda)},
 \qquad Q_a(\lambda)=\int y^2\dd\mu_{\lambda,a}
                  =\sum_j a(x_j(\lambda))y_j(\lambda)^2.
\end{equation}
Here \(z_j=x_j+\ii y_j\) runs over the complete canonical zero multiset.
This is not the Gaussian Gram or pair-correlation statistic of
Section~\ref{sec:localized}; pair interactions enter its time derivative.
A conjugate pair at a point where \(a=1\) contributes exactly \(2y_j^2\),
without background cancellation. The earlier Gram route remains an open
alternative with its stated background and arithmetic-comparison costs.
The results below come from the complete-source calculus and its
assembled internal assessment \cite{SmoothLocalizationProof,SmoothLocalizationAssessment}.
Throughout this section \(0\le\lambda\le1/2\), so the canonical strip
bound is \(\abs{y_j}\le1\) \cite[Theorem 3.2]{Polymath2019}.

\subsection{The exact pair kernel and the complete sum}
Use reverse time \(s=\lambda_*-\lambda\), and put \(f(x,y)=a(x)y^2\).
For distinct roots set \(d=x_j-x_k\), \(e=y_j-y_k\), \(D=d^2+e^2\),
and let \(z_t=z_k+t(z_j-z_k)\). The selected-simple-root dynamics
of \cite[Proposition 3.1]{Polymath2019}, with all other multiplicities
retained, gives
\begin{align}
 \frac{\dd Q_a}{\dd s}&=2\sum_{j<k}'K_f(z_j,z_k),
       \label{eq:full-transverse-kernel}\\
 K_f(z_j,z_k)
 &=\frac{(f_y(z_j)-f_y(z_k))e-(f_x(z_j)-f_x(z_k))d}{D}\nonumber\\
 &=\int_0^1\frac{2a(x_t)e^2-a''(x_t)y_t^2d^2}{D}\dd t.
       \label{eq:transverse-chord}
\end{align}
Indeed differentiation of the endpoint gradients along the chord
produces two mixed terms \(f_{xy}de\) with opposite signs. In a strip
\(\abs y\le Y\),
\begin{equation}\label{eq:gap-free-kernel}
 \abs{K_f}\le\max\{2\norm a_\infty,Y^2\norm{a''}_\infty\}.
\end{equation}
There is no small-gap denominator in this estimate; a direction-independent
value at a coincident pair is neither asserted nor needed.

The prime in \eqref{eq:full-transverse-kernel} is material. Start with
radial cutoffs of the complete opposite-paired Hadamard product and
symmetrize the interaction rows, retaining both members of every far
opposite pair. The complete source bound and Jensen's formula give
\(N_\lambda(r)\le C_*r\log(r+2)\) uniformly on the time interval.
For \(\abs z\le Z\) and \(U\ge2Z+2\), the omitted velocity is at most
\begin{equation}\label{eq:paired-velocity-tail}
 \frac{16C_*Z}{3}\frac{\log(U+2)+1}{U}.
\end{equation}
To see its origin, group
\((z-\zeta)^{-1}+(z+\zeta)^{-1}=2z/(z^2-\zeta^2)\),
then sum the resulting inverse-square majorant by partial summation.
Multiplying \eqref{eq:paired-velocity-tail} by
\(N_\lambda(Z)(Y^2\norm{a'}_\infty+2Y\norm a_\infty)\)
pays the interaction remainder. It is uniform in time at fixed support.
Separately bounding the two chord terms at large separation would lose
this cancellation and can introduce a divergent harmonic sum. Pairs of
inactive endpoints have zero kernel as a whole, even when their chord
passes through the support. These conventions justify the full-index
identity on the simple relevant intervals.

A positive transverse production term does not imply monotonicity of a
localized moment. An exact even heat control is
\[
 P(s,z)=z^4+(12s-15/2)z^2+289/16-15s+12s^2,
 \qquad P_s=P_{zz}.
\]
Its roots at zero are \(\pm2\pm\ii/2\) and remain simple in the strict
strip for \(\abs s\le1/100\). Set
\(b(t)=(1-t)^3(1+t)^{12}\) on \((-1,1)\), zero outside, and
\(a(x)=b(x-2)+b(-x-2)\). This is a nonnegative \(C_c^2\) weight.
Direct substitution gives
\[
 Q_a'(0)=-297/34+288/34=-9/34.
\]
The two terms are horizontal flux and transverse production. This is a
scoped counterexample to generic localized monotonicity, not a canonical
theta orbit.

\subsection{A uniform canonical modulus}
Polymath's unit count and the safe-line estimate in the proof of
Theorem 1.5(iv), Section 9 equation (105), supply a common absolute
\(A\ge1\) such that
\begin{equation}\label{eq:coarse-canonical}
 \abs{H_\lambda'/H_\lambda(u+4\ii)}\le A\log(2+\abs u),\qquad
 \#\{j:u\le x_j\le u+1\}\le A\log(2+\abs u).
\end{equation}
These estimates have no eventual-realness height restriction. The
positive-time constant extends to zero by local uniform convergence of
\(H_\lambda,H_\lambda'\) on the zero-free safe line. Zero counts pass
through enclosing zero-free-boundary rectangles by Rouch\'e's theorem;
a fixed enlargement of \(A\) handles closed endpoints. No convergence
uniformly along an unbounded line is presumed.

Fix \(0\le\chi\le1\), \(\chi\in C_c^2([-1,1])\), and write
\[
 a(x)=\chi((x-X)/R),\quad Q_{X,R}=Q_a,\quad
 \norm\chi_{(2)}=\norm\chi_\infty+\norm{\chi'}_\infty+
                   \norm{\chi''}_\infty.
\]
\begin{theorem}[Canonical modulus through collisions]\label{thm:transverse-lipschitz}
For every \(X\ge10\), \(0<R\le1\) and \(\alpha,\beta\in[0,1/2]\),
with every zero counted with multiplicity,
\begin{equation}\label{eq:transverse-lipschitz}
 \abs{Q_{X,R}(\beta)-Q_{X,R}(\alpha)}
 \le344A^2\norm\chi_{(2)}R^{-2}\log^2(2+X)\abs{\beta-\alpha}.
\end{equation}
The complete interaction formula holds almost everywhere and after
integration. Its first distributional derivative has no collision atom.
\end{theorem}
\begin{proof}
Put \(L_X=\log(2+X)\), and use the bookkeeping cluster
\(\mathcal C_X=\{k:\abs{x_k-X}\le3\}\). Unit covers give population
\(n\le12AL_X\) and active population \(m\le2AL_X\).
For active \(z=x+\ii y\) put \(w=x+4\ii\). The exterior field is
\begin{equation}\label{eq:safe-line-exterior}
 E_X(z)=\frac{H_\lambda'}{H_\lambda}(w)
       -\sum_{k\in\mathcal C_X}\frac1{w-z_k}
       +\sum_{k\notin\mathcal C_X}
          \left(\frac1{z-z_k}-\frac1{w-z_k}\right).
\end{equation}
The last series is absolute. Its terms are bounded by
\(5/\abs{x-x_k}^2\), and unit shells give
\(\sum_{k\notin\mathcal C_X}\abs{x-x_k}^{-2}\le5AL_X\).
The safe-line and finite removal costs are \(2AL_X\) and \(4AL_X\),
so \(\abs{E_X}\le31AL_X\). Truncating the difference series at
horizontal distance \(M\ge2\) leaves at most
\[
 \frac{10A}{M-1}\{L_X+\log(M+1)+1\}.
\]
The paired full field and this absolute remainder justify the grouping;
no absolute harmonic exterior sum is used.

On a simple relevant interval, split the chain rule into finite
interactions and this exterior field. If \(A_i=\norm{a^{(i)}}_\infty\),
then at most \(mn\) local pairs contribute, and
\(\abs{\nabla f}\le A_1+2A_0\). Thus
\[
 \abs{Q_a'}\le A^2L_X^2
 \{48\max(2A_0,A_2)+124(A_1+2A_0)\}
 \le A^2L_X^2(344A_0+124A_1+48A_2).
\]
This implies the stated derivative bound at every such time.

For the assembly, root-multiset continuity in zero-free-boundary
rectangles makes \(Q_a\) continuous, including support crossings.
Repeated-root points in a fixed padded space/time box form a compact
set. At a point of multiplicity \(m\), Polymath's local Hermite
splitting exhausts the cluster with distinct simple roots at every
sufficiently close nonzero time. Every repeated point is therefore
isolated in space and time, so the compact set is finite. Event radii
and Hermite errors need not be uniform. Summing the same derivative
bound on the complementary intervals costs their total duration, not
the number of collisions, and proves \eqref{eq:transverse-lipschitz}.
The essentially uniform remainder \eqref{eq:paired-velocity-tail}
permits cutoff/integration interchange. The first distributional derivative of a Lipschitz function is
represented by an essentially bounded function, completing the claim.
\end{proof}

The microscopic order in \eqref{eq:transverse-lipschitz} is sharp for
an instantaneous coarse-premise class. Take
\(\chi(t)=(1-t^2)^3_+\), \(t_j=7/10+j/[10(m+1)]\), and the even
polynomial with roots \(\pm(X+Rt_j)\pm\ii/2\), where
\(m\ge2\), \(X\ge10\), \(R\le1/10\). On the occupied interval
\(\chi''\ge3\). At its heat snapshot, the same-height, opposite-height
and mirrored contributions respectively give
\[
 Q_s\le-\frac{3m(m-1)}{2R^2}+4m^2+
             \frac{6m^2}{RX}+\frac{2m^2}{X^2}
       \le-\frac{3249m^2}{5000R^2}.
\]
For \(m=\lfloor\log(2+X)\rfloor\), the magnitude is at least
\(\log^2(2+X)/(8R^2)\). Strip, upper counts and safe-line bounds have
common constants at these snapshots. A local heat orbit exists, but its
interval may depend on the parameters; this does not prove sharpness
under canonical hypotheses on the entire fixed time interval.

\subsection{The exact kink and a threshold consumer}
The preceding continuity does not assert differentiability at a
collision. For a multiplicity-\(m\) cluster at \(z_0\), write
\(\tau=\lambda-\lambda_0\). Its leading splitting is
\(z_j=z_0+\sqrt{2\tau}\,h_j+O(\abs\tau)\), where \(h_j\)
are the real probabilists' Hermite roots and
\(\sum h_j^2=m(m-1)\). The contour sum of displacements is analytic:
\(\sum(z_j-z_0)=c\tau+O(\tau^2)\). For any fixed real \(C^2\) test
\(f\), set \(B_f=f_x\Re c+f_y\Im c\), with derivatives at \(z_0\).
Taylor expansion of the cluster moment \(q=\sum f(z_j)\) gives
\begin{equation}\label{eq:collision-kink}
 q'_+=B_f+m(m-1)f_{xx},\qquad
 q'_-=B_f-m(m-1)f_{yy},\qquad
 q'_+-q'_-=m(m-1)\Delta f.
\end{equation}
The analytic displacement sum supplies the common linear term; one
must not differentiate the individual unknown \(O(\abs\tau)\)
remainders. Analytic branches in \(\sqrt\tau\) also justify the
adjacent derivative limits. For \(f=a(x)y^2\) at a real collision,
these slopes are \(-2m(m-1)a(x_0)\) and zero. Simultaneous clusters
add their jumps; endpoints use only the inward slope.

\begin{corollary}[An off-line zero forces a lower threshold]\label{cor:threshold-lower}
If \(H_0\) has a pair \(X\pm\ii d\), \(X\ge10\), \(0<d<1\),
and \(\chi(0)=1\), then
\begin{equation}\label{eq:threshold-lower}
 \Lambda\ge
 \frac{2d^2}{344A^2\norm\chi_{(2)}\log^2(2+X)}.
\end{equation}
\end{corollary}
\begin{proof}
The strip theorem gives \(\Lambda\le1/2\). A nonreal zero at zero
persists for a small positive time, so the hypothesis gives
\(\Lambda>0\). The threshold slice is all real: a nonreal zero there
would persist above it, contradicting the all-real property for
\(\lambda>\Lambda\). Hence \(Q_{X,1}(\Lambda)=0\), whereas
\(Q_{X,1}(0)\ge2d^2\). Apply Theorem~\ref{thm:transverse-lipschitz}.
\end{proof}
This uses the whole fixed-window moment, with no surviving pair label,
maximum-root assumption or isolating tube. For a zeta zero
\(\beta+\ii\gamma\), \(X=2\gamma\), \(d=2\abs{\beta-1/2}\),
so the numerator is \(8\abs{\beta-1/2}^2\). The absolute constant
\(A\) has not been numerically evaluated. This is a conditional lower
bound for \(\Lambda\), not an upper collision-time envelope or an RH
exclusion.
